\documentclass{article}
\usepackage[utf8]{inputenc}
\usepackage{amsmath,amssymb}
\usepackage[most]{tcolorbox}
\usepackage{geometry}
\geometry{margin=2.5cm}

\newcommand{\step}[1]{\par\noindent\hangindent=3.5em\hangafter=1%
  \makebox[3.5em][l]{\textbf{#1.}}}
\newcommand{\steptag}[1]{\hfill{\small\textsf{[#1]}}}

\newtcolorbox{proofbox}[1]{
  colback=gray!5, colframe=gray!60,
  fonttitle=\bfseries, title={#1},
  breakable, enhanced,
  left=4pt, right=4pt, top=4pt, bottom=4pt,
  before skip=10pt, after skip=10pt
}

\begin{document}

\begin{proofbox}{Explicit group-input polar closeness: there exist univers...}
\step{1} Explicit group-input polar closeness: there exist universal C\_grp, C\_pol >= 1, kappa\_pol in (0, 1/2] such that for every finite-dimensional exact-unit epsilon\_r-C*-algebra and delta > 0 satisfying C\_pol*(epsilon\_r + delta) <= kappa\_pol and C\_grp*epsilon\_r < delta - C\_pol*(epsilon\_r*delta + delta\textasciicircum{}2), writing (u\_delta, h\_delta) for the unique inverse of Pi\_delta: calU x B\textasciicircum{}\{calH\}\_delta(J) -> S\_delta := Pi\_delta(calU x B\textasciicircum{}\{calH\}\_delta(J)), Pi\_delta(U, H) = U bold-dot H, the first inverse component u\_delta:S\_delta -> calU satisfies ||u\_delta(U bold-dot V) - U bold-dot V|| <= C\_grp*epsilon\_r and ||u\_delta(U\textasciicircum{}dagger) - U\textasciicircum{}dagger|| <= C\_grp*epsilon\_r for every U, V in calU. \steptag{validated} \\[6pt]
\step{1.1} Constant choice and typed polar inverse: take C\_pol\textasciicircum{}0,kappa\_pol\textasciicircum{}0 from the validated external lem-stage1-polar-retraction, set C\_pol=C\_pol\textasciicircum{}0, kappa\_pol=min(kappa\_pol\textasciicircum{}0,1/16), C\_grp=5, and rho=delta-C\_pol*(epsilon\_r*delta+delta\textasciicircum{}2). These are universal with C\_grp,C\_pol>=1 and kappa\_pol in (0,1/2]. Under the root hypotheses, epsilon\_r+delta<=1/16 and rho>5*epsilon\_r>=0. The hypotheses of lem-stage1-polar-retraction hold, so its displayed map Pi\_delta:calU x B\textasciicircum{}\{calH\}\_delta(J)->S\_delta, Pi\_delta(U,H)=U bold-dot H, has a unique inverse pair (u\_delta,h\_delta):S\_delta->calU x B\textasciicircum{}\{calH\}\_delta(J); for every X in S\_delta this typed pair obeys X=u\_delta(X) bold-dot h\_delta(X), u\_delta(X) in calU, h\_delta(X) in B\textasciicircum{}\{calH\}\_delta(J), hence ||h\_delta(X)-J||<delta, and calU\_rho is contained in S\_delta. \steptag{validated} \\[4pt]
\step{1.2} Raw input estimates and domain typing: with the constants, rho, and typed inverse pair fixed in 1.1, every U,V in calU satisfy ||(U bold-dot V)\textasciicircum{}dagger bold-dot (U bold-dot V)-J||<=3*epsilon\_r and U bold-dot V has a right inverse, while ||(U\textasciicircum{}dagger)\textasciicircum{}dagger bold-dot U\textasciicircum{}dagger-J||<=2*epsilon\_r and U\textasciicircum{}dagger has the right inverse U. Consequently U bold-dot V and U\textasciicircum{}dagger both belong to calU\_rho and hence to S\_delta, so the fixed maps u\_delta,h\_delta may be evaluated at both displayed inputs. \steptag{validated} \\[6pt]
\step{1.2.1} Unitary norm and right-inverse data: for W in calU, def-approximate-unitary-space gives W\textasciicircum{}dagger bold-dot W=J and an R with W bold-dot R=J. The exact-unit clause of def-epsilon-cstar-algebra gives ||J||=1, while its C*-lower bound gives 1=||W\textasciicircum{}dagger bold-dot W||>=(1-epsilon\_r)||W||\textasciicircum{}2; by the smallness derived in 1.1, epsilon\_r<=1/16, so ||W||\textasciicircum{}2<=1/(1-epsilon\_r)<=16/15. \steptag{validated} \\[4pt]
\step{1.2.2} Opposite-product defect: for W in calU and R as in 1.2.1, exact unitality and the associator axiom give ||R-W\textasciicircum{}dagger||=||(W\textasciicircum{}dagger bold-dot W) bold-dot R-W\textasciicircum{}dagger bold-dot (W bold-dot R)||<=epsilon\_r||W||\textasciicircum{}2||R||. Thus epsilon\_r||W||\textasciicircum{}2<=1/15 implies ||R||<=15||W||/14. Since J=W bold-dot R, the product-norm axiom and ||W\textasciicircum{}dagger||=||W|| yield ||W bold-dot W\textasciicircum{}dagger-J||=||W bold-dot(W\textasciicircum{}dagger-R)||<=epsilon\_r*(1+epsilon\_r)||W||\textasciicircum{}3||R||<=((17/16)*(15/14)*(256/225))*epsilon\_r<=2*epsilon\_r. \steptag{validated} \\[4pt]
\step{1.2.3} Product defect: for U,V in calU and X=U bold-dot V, the involution axiom gives X\textasciicircum{}dagger=V\textasciicircum{}dagger bold-dot U\textasciicircum{}dagger. Two applications of the associator axiom compare (V\textasciicircum{}dagger bold-dot U\textasciicircum{}dagger) bold-dot (U bold-dot V) first with ((V\textasciicircum{}dagger bold-dot U\textasciicircum{}dagger) bold-dot U) bold-dot V and then with (V\textasciicircum{}dagger bold-dot (U\textasciicircum{}dagger bold-dot U)) bold-dot V=(V\textasciicircum{}dagger bold-dot J) bold-dot V=J. The product-norm axiom and 1.2.1 therefore give ||X\textasciicircum{}dagger bold-dot X-J||<=2*epsilon\_r*(1+epsilon\_r)||U||\textasciicircum{}2||V||\textasciicircum{}2<=(544/225)*epsilon\_r<=3*epsilon\_r. \steptag{validated} \\[6pt]
\step{1.2.3.1} Self-contained numerical bridge (independent of sibling 1.2.1): the fixed root guard and constant choice give C\_pol>=1 and C\_pol*(epsilon\_r+delta)<=kappa\_pol<=1/16, hence 0<=epsilon\_r<=1/16. For each W in calU, the definition of calU gives W\textasciicircum{}dagger bold-dot W=J, while exact unitality and the C*-lower axiom give 1=||J||=||W\textasciicircum{}dagger bold-dot W||>=(1-epsilon\_r)||W||\textasciicircum{}2; since 1-epsilon\_r>0, ||W||\textasciicircum{}2<=1/(1-epsilon\_r)<=16/15. Thus for U,V in calU both ||U||\textasciicircum{}2 and ||V||\textasciicircum{}2 are at most 16/15. Now set A=(V\textasciicircum{}dagger bold-dot U\textasciicircum{}dagger) bold-dot (U bold-dot V), B=((V\textasciicircum{}dagger bold-dot U\textasciicircum{}dagger) bold-dot U) bold-dot V, and C=(V\textasciicircum{}dagger bold-dot (U\textasciicircum{}dagger bold-dot U)) bold-dot V. The associator and product-norm axioms give ||A-B||<=epsilon\_r||V\textasciicircum{}dagger bold-dot U\textasciicircum{}dagger||||U||||V||<=epsilon\_r*(1+epsilon\_r)||U||\textasciicircum{}2||V||\textasciicircum{}2 and ||B-C||<=(1+epsilon\_r)||((V\textasciicircum{}dagger bold-dot U\textasciicircum{}dagger) bold-dot U)-V\textasciicircum{}dagger bold-dot(U\textasciicircum{}dagger bold-dot U)|| ||V||<=epsilon\_r*(1+epsilon\_r)||U||\textasciicircum{}2||V||\textasciicircum{}2. Since U\textasciicircum{}dagger bold-dot U=J and exact unitality gives C=(V\textasciicircum{}dagger bold-dot J) bold-dot V=V\textasciicircum{}dagger bold-dot V=J, while the involution axiom gives A=(U bold-dot V)\textasciicircum{}dagger bold-dot(U bold-dot V), the triangle inequality yields the defect bound <=2*epsilon\_r*(1+epsilon\_r)*(16/15)\textasciicircum{}2<=(544/225)*epsilon\_r<=3*epsilon\_r. This proves the numerical conclusion entirely inside the subtree of 1.2.3, including when epsilon\_r=0. \steptag{validated} \\[4pt]
\step{1.2.4} Exact right inverse for the product: for X=U bold-dot V, the product-norm estimate and 1.2.1 give ||X||\textasciicircum{}2<=(1+epsilon\_r)\textasciicircum{}2*(16/15)\textasciicircum{}2<=289/225. If X bold-dot Y=0, exact unitality, the product-norm axiom, the associator axiom, and 1.2.3 imply ||Y||<=||(J-X\textasciicircum{}dagger bold-dot X) bold-dot Y||+||(X\textasciicircum{}dagger bold-dot X) bold-dot Y-X\textasciicircum{}dagger bold-dot (X bold-dot Y)||<=((1+epsilon\_r)*3*epsilon\_r+epsilon\_r||X||\textasciicircum{}2)||Y||<=(51/256+289/3600)||Y||<(1/2)||Y||. Hence Y=0. Bilinearity makes L\_X:Y->X bold-dot Y linear; because the algebra is finite-dimensional, injectivity makes L\_X surjective, so some R\_X satisfies X bold-dot R\_X=J. \steptag{validated} \\[4pt]
\step{1.2.5} Domain conclusion for the fixed inverse: X\_*=U\textasciicircum{}dagger has the exact right inverse U, and the involution axiom plus 1.2.2 gives ||(X\_*)\textasciicircum{}dagger bold-dot X\_*-J||=||U bold-dot U\textasciicircum{}dagger-J||<=2*epsilon\_r. The product X\_p=U bold-dot V has a right inverse by 1.2.4 and defect at most 3*epsilon\_r by 1.2.3. Since rho>5*epsilon\_r by 1.1, each defect is strictly below 2*rho (when epsilon\_r=0, both defects vanish and rho>0). By def-approximate-unitary-space both inputs lie in calU\_rho; the inclusion from lem-stage1-polar-retraction recorded in 1.1 puts both in S\_delta, where the already-fixed typed pair (u\_delta,h\_delta) is defined. \steptag{validated} \\[6pt]
\step{1.2.5.1} Validated-premise domain bridge: fix U,V in calU and put X\_p=U bold-dot V and X\_*=U\textasciicircum{}dagger. By validated 1.2.2 instantiated at W=U, ||U bold-dot U\textasciicircum{}dagger-J||<=2*epsilon\_r. The involution axiom gives (X\_*)\textasciicircum{}dagger=U, while U in calU gives U\textasciicircum{}dagger bold-dot U=J; hence ||(X\_*)\textasciicircum{}dagger bold-dot X\_*-J||<=2*epsilon\_r and X\_* bold-dot U=J, so U is a right inverse of X\_*. By validated 1.2.3 and 1.2.4, ||X\_p\textasciicircum{}dagger bold-dot X\_p-J||<=3*epsilon\_r and X\_p has a right inverse. Validated 1.1 gives rho>5*epsilon\_r and calU\_rho subseteq S\_delta for its fixed typed inverse pair. If epsilon\_r>0, both 2*epsilon\_r and 3*epsilon\_r are strictly less than 2*rho; if epsilon\_r=0, both defects vanish and rho>0, so the same strict inequalities hold. Therefore def-approximate-unitary-space places X\_* and X\_p in calU\_rho, hence in S\_delta, and the single typed pair (u\_delta,h\_delta) fixed in 1.1 is defined at both inputs. \steptag{validated} \\[6pt]
\step{1.2.5.1.1} Replacement product-input estimate independent of pending 1.2.3: put e=epsilon\_r and X=U bold-dot V. Validated 1.1 gives 0<=e<=1/16. Since U,V are in calU, def-approximate-unitary-space gives U\textasciicircum{}dagger bold-dot U=V\textasciicircum{}dagger bold-dot V=J; the C*-lower axiom and ||J||=1 give ||U||\textasciicircum{}2,||V||\textasciicircum{}2<=1/(1-e). The involution axiom gives X\textasciicircum{}dagger=V\textasciicircum{}dagger bold-dot U\textasciicircum{}dagger. Applying the associator axiom first to (V\textasciicircum{}dagger bold-dot U\textasciicircum{}dagger,U,V), and then to (V\textasciicircum{}dagger,U\textasciicircum{}dagger,U) followed by the product-norm axiom for right multiplication by V, yields ||X\textasciicircum{}dagger bold-dot X-J||<=2*e*(1+e)*||U||\textasciicircum{}2*||V||\textasciicircum{}2<=2*e*(17/16)*(16/15)\textasciicircum{}2=(544/225)*e<3*e when e>0, while it is 0 when e=0. Moreover ||X||\textasciicircum{}2<=(1+e)\textasciicircum{}2||U||\textasciicircum{}2||V||\textasciicircum{}2<=(17/15)\textasciicircum{}2=289/225. If X bold-dot Y=0, exact unitality, the product-norm axiom and the associator axiom give ||Y||<=((1+e)*3*e+e||X||\textasciicircum{}2)||Y||<=(51/256+289/3600)||Y||<(1/2)||Y||, hence Y=0. Thus the finite-dimensional linear map L\_X is injective and therefore surjective, so X has a right inverse. This proves both product-input requirements without 1.2.3 or 1.2.4. \steptag{validated} \\[4pt]
\step{1.2.5.1.2} Corrected domain conclusion using only validated premises and the replacement child: for X\_*=U\textasciicircum{}dagger, validated 1.2.2 gives ||(X\_*)\textasciicircum{}dagger bold-dot X\_*-J||=||U bold-dot U\textasciicircum{}dagger-J||<=2*epsilon\_r, and X\_* bold-dot U=U\textasciicircum{}dagger bold-dot U=J gives a right inverse. The preceding replacement child gives X\_p=U bold-dot V a right inverse and defect at most 3*epsilon\_r. Validated 1.1 gives rho>5*epsilon\_r and calU\_rho subseteq S\_delta. Hence both defects are strictly below 2*rho (also when epsilon\_r=0, since then they vanish and rho>0), so def-approximate-unitary-space places X\_* and X\_p in calU\_rho and therefore in S\_delta, where the single typed inverse pair fixed by 1.1 is defined. No use is made of pending node 1.2.3. \steptag{validated} \\[4pt]
\step{1.3} Polar closeness and conclusion: under epsilon\_r+delta<=1/16, if X lies in S\_delta and ||X\textasciicircum{}dagger bold-dot X-J||<=3*epsilon\_r, then the typed inverse pair fixed in 1.1 satisfies ||u\_delta(X)-X||<=5*epsilon\_r. Apply this first to X=U bold-dot V and then to X=U\textasciicircum{}dagger using 1.2; since C\_grp=5, both inequalities required by the root follow for the same displayed u\_delta. \steptag{validated} \\[6pt]
\step{1.3.1} Polar square comparison for the fixed pair: write u=u\_delta(X), h=h\_delta(X), and a=h-J for the typed inverse pair supplied in 1.1. Then X=u bold-dot h, u lies in calU, h is Hermitian, and ||a||<delta. Thus ||u||\textasciicircum{}2<=1/(1-epsilon\_r)<=16/15 by the C*-lower bound applied to u\textasciicircum{}dagger bold-dot u=J, while ||h||<=1+delta<=17/16. The involution axiom gives X\textasciicircum{}dagger=h bold-dot u\textasciicircum{}dagger. Two associator comparisons, from (h bold-dot u\textasciicircum{}dagger) bold-dot (u bold-dot h) to ((h bold-dot u\textasciicircum{}dagger) bold-dot u) bold-dot h and then to (h bold-dot (u\textasciicircum{}dagger bold-dot u)) bold-dot h=h bold-dot h, together with the product-norm axiom give ||X\textasciicircum{}dagger bold-dot X-h bold-dot h||<=2*epsilon\_r*(1+epsilon\_r)||h||\textasciicircum{}2||u||\textasciicircum{}2<=(9826/3840)*epsilon\_r<=3*epsilon\_r. \steptag{validated} \\[4pt]
\step{1.3.2} Quadratic absorption: exact unitality and bilinearity give h bold-dot h-J=(J+a) bold-dot (J+a)-J=2a+a bold-dot a. Combining the assumed defect ||X\textasciicircum{}dagger bold-dot X-J||<=3*epsilon\_r with 1.3.1 yields 2||a||<=6*epsilon\_r+(1+epsilon\_r)||a||\textasciicircum{}2. The guard in 1.1 gives ||a||<delta<=1/16 and epsilon\_r<=1/16, hence (1+epsilon\_r)||a||\textasciicircum{}2<=(17/256)||a||. Therefore (495/256)||a||<=6*epsilon\_r and ||a||<=(1536/495)*epsilon\_r<=4*epsilon\_r. \steptag{validated} \\[4pt]
\step{1.3.3} Return to the same first inverse component: exact right unitality and the typed identity X=u\_delta(X) bold-dot h\_delta(X) from 1.1 give ||u\_delta(X)-X||=||u bold-dot J-u bold-dot h||=||u bold-dot(J-h)||<=(1+epsilon\_r)||u||||a||. Using ||u||<=sqrt(16/15) from 1.3.1, epsilon\_r<=1/16 from 1.1, and 1.3.2 gives ||u\_delta(X)-X||<=(17/sqrt(15))*epsilon\_r<=5*epsilon\_r. Applying this bound to the two S\_delta inputs typed in 1.2 (whose defects are at most 3*epsilon\_r) proves the two root inequalities for the single inverse component u\_delta fixed in 1.1. \steptag{validated} \\[4pt]
\end{proofbox}

\end{document}
